\section{Discarded directions and limitations}
\label{app:alternatives}
\subsection{Total-variation budgets and reconstruction operators}
While looking for other ways of controlling information, we asked whether an image could be simplified while staying as similar as possible to the original. One family keeps each channel mean and bounds the total variation,
\[
 \mathcal K_t(x)=\left\{z\in[0,1]^{3\times H\times W}:
 \operatorname{TV}(z)\le t\operatorname{TV}(x),\ \bar z_c=\bar x_c\right\},
\]
and returns the element of $\mathcal K_t(x)$ closest to $x$, either in the $L^2$ norm or in the homogeneous $\dot H^{-1}$ norm $\frac12\sum_c(z_c-x_c)^\top\mathsf L^\dagger(z_c-x_c)$, where $\mathsf L$ is the forward-difference Laplacian. We only computed previews of these transformations on ten images. Each solve needed about $10^4$ primal--dual iterations, which is too slow for training, and no network was trained with them.

The $L^2$ and $\dot H^{-1}$ \emph{reconstruction operators} of Appendix~\ref{app:explore-operators} are a different experiment. They are resolution reductions used inside the network during training, with no total-variation budget: for each channel of a map $h$, they return the $r\times r$ coefficients $z$ minimizing $\frac12(U_rz-h)^\top Q\,(U_rz-h)$ under $\mathbf 1^\top(U_rz-h)=0$, where $U_r$ is bilinear upsampling and $Q=I$ or $Q=\mathsf L^\dagger$.

\subsection{Other directions}
We tested a db2 wavelet shrinkage at the 19 convolution outputs in one one-seed pilot on a 10,000-image subset. It reached $\xDbWave\%$ validation accuracy against $\xDbPlain\%$ for plain, but a run took $\xDbWaveWall$\,s against $\xDbPlainWall$\,s, and we set it aside. Schedules chosen from training signals and different allocations of time between blur levels, run at seed 0 only, finished between $\xAdMin\%$ and $\xAdMax\%$; three reruns of the fixed plateau schedule in the same study finished between $\xCalMin\%$ and $\xCalMax\%$, so these runs do not show a clear improvement. In the ablation batch, per-site scale profiles without reduction ($\xAbQMin$--$\xAbQMax\%$) stayed below uniform post-ReLU filtering. Early investigations on other branches, acting on residual-branch amplitudes, activations, anchoring to the initialization or example ordering, helped narrow the focus; their records are not part of our audited results and we do not report them.

\subsection{Limitations}
\textcolor{red}{a diluer une fois qu'on aura fait des benchmarks bien solides. Cette partie de l'annexe concerne la procédure qui a servi à retenir des méthodes exclusivement. }
All results use one dataset and one architecture, without augmentation, and a short budget: the plain baseline stays far below a tuned ResNet-20. The gains measured after 30 epochs need not carry over to longer training. In one-seed runs of 120 epochs, the plateau Gaussian at the convolution outputs reached $\xLongLoPlateau\%$ against $\xLongLoPlain\%$ for plain at the reference learning rate, and $\xLongHiPlateau\%$ against $\xLongHiPlain\%$ with a learning rate ten times larger. We used at most three seeds, results vary between batches, and the test set was used for selection. BatchNorm statistics accumulate under each intervention, which complicates any comparison before the final epoch. Finally, the reduced-resolution phases did not translate into shorter training on our T4 GPUs, and the combined method took longer than plain ($\uWallComb$\,s against $\uWallPlain$\,s per run, evaluation included).
